% =====================================================================
\chapter{Soluciones de los ejercicios}
\label{ap:soluciones}
% =====================================================================

Cada solución se presenta en cinco partes: datos, planteamiento, desarrollo, resultado y comprobación. En los pasos intermedios se redondea a cuatro cifras significativas salvo que se indique otra cosa; los resultados exactos se dan como fracción o como expresión cerrada cuando existe. Todos los valores numéricos se han comprobado con un programa en Python.

% ---------------------------------------------------------------------
\section{Tema 1}

\begin{solucion}{eje:1-kv}
\datos $o=\sum_j\alpha_jv_j$ con $\alpha_j=\softmax_j\bigl(a(q,k_j)\bigr)$; $k_j=x_j\WK$, $v_j=x_j\WV$.
\planteamiento localizar en qué parte de la \cref{def:agrupacion} interviene cada vector.
\pasos
\begin{enumerate}[label=\arabic*.]
  \item La clave $k_j$ solo aparece en la puntuación $a(q,k_j)$: determina \emph{si} la posición $j$ recibe peso y cuánto.
  \item El valor $v_j$ solo aparece en la suma ponderada: determina \emph{qué} información aporta la posición $j$ cuando recibe peso.
  \item Con $\WK=\WV$, ambos vectores coinciden: las características por las que un token es encontrado serían las mismas que transmite. Un token no podría, por ejemplo, anunciarse por su categoría gramatical y transmitir su significado.
\end{enumerate}
\resultado \textbf{la clave fija el peso y el valor fija el contenido}; imponer $\WK=\WV$ obliga a que el criterio de búsqueda y la información transmitida vivan en el mismo subespacio, lo que reduce la capacidad del modelo.
\comprobacion en el \cref{ej:tres}, la tercera clave $(1,1,1,1)$ hace que la posición~3 sea «medio parecida» a todas las consultas, mientras que su valor $(1,1)$ aporta las dos coordenadas: los papeles son distintos.
\end{solucion}

\begin{solucion}{eje:1-escalar}
\datos $q=(1,1)$; $k_1=(1,0)$, $k_2=(0,1)$, $k_3=(1,1)$; valores $10$, $20$, $30$; $a(q,k)=q\cdot k/\sqrt2$.
\planteamiento puntuaciones, softmax y media ponderada (\cref{def:agrupacion}).
\pasos
\begin{enumerate}[label=\arabic*.]
  \item Productos escalares: $q\cdot k_1=1$, $q\cdot k_2=1$, $q\cdot k_3=2$.
  \item Escalado: $1/\sqrt2=0{,}7071$ (dos veces) y $2/\sqrt2=1{,}414$.
  \item Exponenciales: $e^{0{,}7071}=2{,}028$ (dos veces) y $e^{1{,}414}=4{,}113$; suma $2{,}028+2{,}028+4{,}113=8{,}169$.
  \item Pesos: $2{,}028/8{,}169=0{,}2483$ (dos veces) y $4{,}113/8{,}169=0{,}5035$.
  \item Salida: $0{,}2483\cdot 10+0{,}2483\cdot 20+0{,}5035\cdot 30=2{,}483+4{,}966+15{,}10=22{,}55$.
\end{enumerate}
\resultado pesos \textbf{$(0{,}2483;\allowbreak\ 0{,}2483;\allowbreak\ 0{,}5035)$}; salida \textbf{$22{,}55$}.
\comprobacion los pesos suman $1{,}000$ y la salida está entre $10$ y $30$ (\cref{prop:convexa}); supera la media simple, $20$, porque $k_3$ es la clave más parecida a $q$.
\end{solucion}

\begin{solucion}{eje:1-aditiva}
\datos $W_a=U_a=I_2$, $v_a=(1,1)\T$, $s_{i-1}=(1,-1)\T$; $h_1=(1,0)\T$, $h_2=(0,1)\T$, $h_3=(1,1)\T$.
\planteamiento $e_{ij}=v_a\T\tanh(s_{i-1}+h_j)$ es la suma de las dos componentes de $\tanh(s_{i-1}+h_j)$; después, softmax y media ponderada (\cref{def:aditiva}).
\pasos
\begin{enumerate}[label=\arabic*.]
  \item $s_{i-1}+h_1=(2,-1)$, $\tanh=(0{,}9640;\allowbreak\ -0{,}7616)$, $e_{i1}=0{,}2024$.
  \item $s_{i-1}+h_2=(1,0)$, $\tanh=(0{,}7616;\allowbreak\ 0)$, $e_{i2}=0{,}7616$.
  \item $s_{i-1}+h_3=(2,0)$, $\tanh=(0{,}9640;\allowbreak\ 0)$, $e_{i3}=0{,}9640$.
  \item Exponenciales: $e^{0{,}2024}=1{,}224$, $e^{0{,}7616}=2{,}142$, $e^{0{,}9640}=2{,}622$; suma $5{,}988$.
  \item Pesos: $1{,}224/5{,}988=0{,}2045$, $2{,}142/5{,}988=0{,}3577$, $2{,}622/5{,}988=0{,}4379$.
  \item Contexto: $c_i=0{,}2045\,(1,0)+0{,}3577\,(0,1)+0{,}4379\,(1,1)=(0{,}2045+0{,}4379;\allowbreak\ 0{,}3577+0{,}4379)=(0{,}6424;\allowbreak\ 0{,}7955)$.
\end{enumerate}
\resultado $e_{i\cdot}=(0{,}2024;\allowbreak\ 0{,}7616;\allowbreak\ 0{,}9640)$, $\alpha_{i\cdot}=(0{,}2045;\allowbreak\ 0{,}3577;\allowbreak\ 0{,}4379)$ y \textbf{$c_i=(0{,}6424;\allowbreak\ 0{,}7955)$}.
\comprobacion los pesos suman $1{,}000$, y cada coordenada de $c_i$ está en $[0,1]$, el rango de las coordenadas de los $h_j$ (\cref{prop:convexa}).
\end{solucion}

\begin{solucion}{eje:1-copiar}
\datos $k_1,k_2,k_3$ base canónica de $\R^{3}$; $v_1=(2,0)$, $v_2=(0,4)$, $v_3=(6,6)$; $\alpha=\softmax(q\cdot k_1,q\cdot k_2,q\cdot k_3)$.
\planteamiento con claves ortonormales, $q\cdot k_j$ es la coordenada $j$ de $q$; la concentración del peso depende de las diferencias entre puntuaciones (\cref{teo:limites}).
\pasos
\begin{enumerate}[label=\arabic*.]
  \item (a) Condición: $q\cdot k_j\gg q\cdot k_l$ para todo $l\ne j$; es decir, $q$ alineada con $k_j$, de norma grande y casi ortogonal a las demás claves.
  \item Con $q=10k_3=(0,0,10)$: puntuaciones $0,0,10$; exponenciales $1,1,e^{10}=22\,026$; suma $22\,028$. Pesos $4{,}540\cdot 10^{-5}$ (dos veces) y $0{,}99991$.
  \item Salida: $4{,}540\cdot10^{-5}(2,0)+4{,}540\cdot10^{-5}(0,4)+0{,}99991\,(6,6)=(5{,}9995;\allowbreak\ 5{,}9996)\approx v_3$.
  \item (b) Con $q=10(k_1+k_2)=(10,10,0)$: exponenciales $e^{10},e^{10},1$, suma $44\,053$; pesos $0{,}49999$ (dos veces) y $2{,}270\cdot10^{-5}$.
  \item Salida: $0{,}49999(2,0)+0{,}49999(0,4)+2{,}270\cdot10^{-5}(6,6)=(1{,}0001;\allowbreak\ 2{,}0001)\approx\tfrac12(v_1+v_2)=(1,2)$.
  \item (c) Con $q=k_1+k_2=(1,1,0)$: exponenciales $e=2{,}718$ (dos veces) y $1$, suma $6{,}437$; pesos $0{,}4223$, $0{,}4223$, $0{,}1554$.
  \item Salida: $0{,}4223(2,0)+0{,}4223(0,4)+0{,}1554(6,6)=(0{,}8446+0{,}9322;\allowbreak\ 1{,}689+0{,}9322)=(1{,}777;\allowbreak\ 2{,}621)$.
\end{enumerate}
\resultado (a) $q$ alineada con una clave y de norma grande; con $10k_3$, \textbf{$c\approx(6,6)=v_3$}. (b) $q=10(k_1+k_2)$ da \textbf{$c\approx(1,2)$}. (c) Sin el factor~10, \textbf{$c=(1{,}777;\allowbreak\ 2{,}621)$}: el valor $v_3$ se filtra con un peso del $15{,}5\,\%$.
\comprobacion en los tres casos los pesos suman~1. La norma de $q$ actúa como el factor $\beta$ del \cref{teo:limites}; en un modelo entrenado, la red la ajusta agrandando $\WQ$.
\end{solucion}

\begin{solucion}{eje:1-limites}
\datos $\alpha_j(\beta)=e^{\beta s_j}/\sum_le^{\beta s_l}$; $s^{\ast}=\max_js_j$, alcanzado en $J^{\ast}$ con $\abs{J^{\ast}}=r$.
\planteamiento repetir la demostración del \cref{teo:limites} sin suponer unicidad.
\pasos
\begin{enumerate}[label=\arabic*.]
  \item Dividiendo por $e^{\beta s^{\ast}}$: $\alpha_j(\beta)=e^{\beta(s_j-s^{\ast})}/\sum_le^{\beta(s_l-s^{\ast})}$.
  \item Si $j\in J^{\ast}$, el numerador vale $e^{0}=1$ para todo $\beta$; si $j\notin J^{\ast}$, $s_j-s^{\ast}<0$ y el numerador tiende a~0 cuando $\beta\to\infty$.
  \item El denominador es la suma de $r$ unos y de términos que tienden a~0: tiende a $r$.
  \item Luego $\alpha_j\to 1/r$ si $j\in J^{\ast}$ y $\alpha_j\to 0$ en otro caso, y la salida tiende a $\sum_{j\in J^{\ast}}\tfrac1r v_j$.
  \item Si $v_1=\dots=v_n=v$, la envolvente convexa de la \cref{prop:convexa} es el punto $\{v\}$, luego la salida vale $v$ para toda consulta.
\end{enumerate}
\resultado $\displaystyle\lim_{\beta\to\infty}\Att_\beta(q)=\frac1r\sum_{j\in J^{\ast}}v_j$; con valores iguales, la salida es constante.
\comprobacion con $r=1$ se recupera el \cref{teo:limites}. En la fila~3 del \cref{ej:tres}, las tres puntuaciones empatan y la salida es la media $(0{,}667;\allowbreak\ 0{,}667)$ para cualquier $\beta$, de acuerdo con el resultado con $r=3$.
\end{solucion}

% ---------------------------------------------------------------------
\section{Tema 2}

\begin{solucion}{eje:2-formas}
\datos $n=10$, $d_k=16$, $d_v=32$; $Q,K\in\R^{10\times16}$, $V\in\R^{10\times32}$.
\planteamiento regla de formas $(a\times b)(b\times c)=a\times c$; la softmax por filas conserva la forma.
\pasos
\begin{enumerate}[label=\arabic*.]
  \item $S=QK\T$: $(10\times16)(16\times10)=10\times10$.
  \item $A=\softmax(S/4)$: $10\times10$.
  \item $O=AV$: $(10\times10)(10\times32)=10\times32$.
\end{enumerate}
\resultado \textbf{$S$ y $A$ son $10\times10$; $O$ es $10\times32$.} La dimensión $d_k$ se contrae en el producto escalar $q_i\cdot k_j$ y no sobrevive en la salida.
\comprobacion $O$ tiene una fila por posición, y cada fila es una combinación de valores, de dimensión $d_v=32$ (\cref{tab:formas}).
\end{solucion}

\begin{solucion}{eje:2-var}
\datos componentes independientes de media~0 y varianza~1; $d_k\in\{64,128\}$.
\planteamiento por el \cref{teo:varianza}, $\Var(q\cdot k)=d_k$ y $\Var(q\cdot k/\sqrt{d_k})=1$.
\pasos
\begin{enumerate}[label=\arabic*.]
  \item $d_k=64$: varianza $64$, desviación típica $\sqrt{64}=8$.
  \item $d_k=128$: varianza $128$, desviación típica $\sqrt{128}=8\sqrt2=11{,}31$.
  \item Tras escalar: varianza $d_k/d_k=1$ y desviación típica~1 en ambos casos.
\end{enumerate}
\resultado \textbf{varianzas $64$ y $128$ (desviaciones $8$ y $11{,}31$); tras escalar, $1$ y $1$.}
\comprobacion la simulación citada tras el \cref{teo:varianza} da $63{,}8$ para $d_k=64$, a menos de un $0{,}5\,\%$ del valor teórico.
\end{solucion}

\begin{solucion}{eje:2-temp}
\datos $z=(1{,}2;\allowbreak\ 0{,}9;\allowbreak\ 0{,}1;\allowbreak\ -0{,}5)$; $\tau\in\{0{,}2;\allowbreak\ 2\}$.
\planteamiento $\softmax(z/\tau)_i=e^{z_i/\tau}/\sum_je^{z_j/\tau}$; dividir por $\tau$ es multiplicar por $\beta=1/\tau$.
\pasos
\begin{enumerate}[label=\arabic*.]
  \item $\tau=0{,}2$: $z/\tau=(6;\allowbreak\ 4{,}5;\allowbreak\ 0{,}5;\allowbreak\ -2{,}5)$. Exponenciales $403{,}4$; $90{,}02$; $1{,}649$; $0{,}08208$; suma $495{,}2$.
  \item Probabilidades: $0{,}8147$; $0{,}1818$; $0{,}003330$; $0{,}0001657$.
  \item $\tau=2$: $z/\tau=(0{,}6;\allowbreak\ 0{,}45;\allowbreak\ 0{,}05;\allowbreak\ -0{,}25)$. Exponenciales $1{,}822$; $1{,}568$; $1{,}051$; $0{,}7788$; suma $5{,}221$.
  \item Probabilidades: $0{,}3490$; $0{,}3004$; $0{,}2014$; $0{,}1492$.
  \item $\tau\to0^{+}$ equivale a $\beta\to\infty$: por el \cref{teo:limites}, la distribución tiende a $(1,0,0,0)$, el \emph{argmax}. $\tau\to\infty$ equivale a $\beta\to0$: tiende a la uniforme $(0{,}25;\allowbreak\ 0{,}25;\allowbreak\ 0{,}25;\allowbreak\ 0{,}25)$.
\end{enumerate}
\resultado \textbf{$\tau=0{,}2$: $(0{,}8147;\allowbreak\ 0{,}1818;\allowbreak\ 0{,}0033;\allowbreak\ 0{,}0002)$; $\tau=2$: $(0{,}3490;\allowbreak\ 0{,}3004;\allowbreak\ 0{,}2014;\allowbreak\ 0{,}1492)$}; los límites son el \emph{argmax} y la uniforme. En la atención, $\sqrt{d_k}$ desempeña el papel de $\tau$: sin él, las puntuaciones crecen con $d_k$ y es como trabajar con una temperatura demasiado baja.
\comprobacion las dos distribuciones suman $1{,}000$; con $\tau=1$ se obtiene $(0{,}4432;\allowbreak\ 0{,}3283;\allowbreak\ 0{,}1475;\allowbreak\ 0{,}0810)$, intermedia entre ambas.
\end{solucion}

\begin{solucion}{eje:2-jac}
\datos $z=(1,0,0)$; $J_{ij}=p_i(\delta_{ij}-p_j)$, es decir, $J=\diag(p)-p\,p\T$.
\planteamiento calcular $p$ y sustituir.
\pasos
\begin{enumerate}[label=\arabic*.]
  \item $p=\bigl(e,1,1\bigr)/(e+2)$: $e+2=4{,}718$, luego $p=(0{,}5761;\allowbreak\ 0{,}2119;\allowbreak\ 0{,}2119)$.
  \item Diagonal: $p_1(1-p_1)=0{,}5761\cdot0{,}4239=0{,}2442$; $p_2(1-p_2)=0{,}2119\cdot0{,}7881=0{,}1670$, igual para $p_3$.
  \item Fuera de la diagonal: $-p_1p_2=-0{,}1221$, $-p_1p_3=-0{,}1221$, $-p_2p_3=-0{,}04492$.
  \item Cota: $1-\max_lp_l=1-0{,}5761=0{,}4239\ge0{,}2442=\max\abs{J_{ij}}$.
  \item Suma de la columna~1: $0{,}2442-0{,}1221-0{,}1221=0$; de la columna~2: $-0{,}1221+0{,}1670-0{,}04492=0$ (con el redondeo).
  \item Motivo: $\sum_ip_i=1$ para todo $z$, luego $\sum_i\partial p_i/\partial z_j=\partial(1)/\partial z_j=0$.
\end{enumerate}
\resultado
\[
J=\begin{pmatrix}0{,}2442&-0{,}1221&-0{,}1221\\-0{,}1221&0{,}1670&-0{,}04492\\-0{,}1221&-0{,}04492&0{,}1670\end{pmatrix},\qquad \max\abs{J_{ij}}=0{,}2442\le0{,}4239 .
\]
\comprobacion el máximo coincide con el de la \cref{tab:saturacion}. La matriz es simétrica, como corresponde a la hessiana de $\LSE$ (módulo~4).
\end{solucion}

\begin{solucion}{eje:2-sigma}
\datos $q_i,k_i$ independientes, de media~0 y varianza $\sigma^{2}$, $i=1,\dots,d_k$.
\planteamiento repetir la demostración del \cref{teo:varianza} con $\sigma^{2}$.
\pasos
\begin{enumerate}[label=\arabic*.]
  \item $\E[q_ik_i]=\E[q_i]\,\E[k_i]=0$.
  \item $\E[q_i^{2}]=\Var(q_i)+\E[q_i]^{2}=\sigma^{2}$, igual para $k_i$.
  \item $\Var(q_ik_i)=\E[q_i^{2}]\,\E[k_i^{2}]-0=\sigma^{4}$.
  \item Los $d_k$ sumandos son independientes: $\Var(q\cdot k)=d_k\sigma^{4}$.
  \item $\Var(c\,q\cdot k)=c^{2}d_k\sigma^{4}=1$ si y solo si $c=1/(\sigma^{2}\sqrt{d_k})$ (tomando $c>0$).
\end{enumerate}
\resultado \textbf{$\Var(q\cdot k)=d_k\sigma^{4}$ y $c=1/(\sigma^{2}\sqrt{d_k})$.}
\comprobacion con $\sigma=1$ se recupera $1/\sqrt{d_k}$. Con $\sigma=2$ y $d_k=16$, la fórmula da $16\cdot16=256$; una simulación con $2\cdot10^{5}$ pares da $256{,}6$. En la práctica se usa $1/\sqrt{d_k}$ y las capas de normalización (módulo~11) mantienen $\sigma\approx1$.
\end{solucion}

\begin{solucion}{eje:2-cruzada}
\datos $Q\in\R^{m\times d_k}$; $K=X\WK$, $V=X\WV$ con $X\in\R^{n\times d}$; $P\in\R^{n\times n}$ de permutación.
\planteamiento sustituir $X$ por $PX$, de modo que $K'=PK$ y $V'=PV$, y seguir cada operación como en la \cref{prop:equivariancia}.
\pasos
\begin{enumerate}[label=\arabic*.]
  \item Puntuaciones: $S'=Q(PK)\T=QK\T P\T=SP\T$. Multiplicar por $P\T$ a la derecha permuta las \emph{columnas} de $S$.
  \item La softmax de cada fila es equivariante respecto de permutaciones de sus entradas: $A'=\softmax(S'/\sqrt{d_k})=AP\T$.
  \item Salida: $O'=A'V'=AP\T PV=AV=O$, porque $P\T P=I$.
\end{enumerate}
\resultado \textbf{$\Att(Q,PK,PV)=\Att(Q,K,V)$}: la atención cruzada trata la memoria $(K,V)$ como un conjunto no ordenado.
\comprobacion en la autoatención la permutación afecta también a las consultas y aparece un factor $P$ a la izquierda: se recupera $O'=PO$, la equivariancia de la \cref{prop:equivariancia}.
\end{solucion}

% ---------------------------------------------------------------------
\section{Tema 3}

\begin{solucion}{eje:3-despues}
\datos puntuaciones escaladas $(0,4,2)$; el tercer token es futuro.
\planteamiento comparar cada versión con la correcta, $\softmax(0,4,-\infty)=(0{,}018;\allowbreak\ 0{,}982;\allowbreak\ 0)$ (\cref{ej:mascara}).
\pasos
\begin{enumerate}[label=\arabic*.]
  \item (a) $\softmax(0,4,2)=(0{,}01588;\allowbreak\ 0{,}8668;\allowbreak\ 0{,}1173)$. Anulando el tercero queda $(0{,}01588;\allowbreak\ 0{,}8668;\allowbreak\ 0)$, que suma $0{,}8827$.
  \item Renormalizando: $0{,}01588/0{,}8827=0{,}01799$ y $0{,}8668/0{,}8827=0{,}9820$, que es el resultado correcto; sin renormalizar, la fila no es una distribución.
  \item (b) Puntuaciones $(0,4,0)$: $\softmax=(0{,}01767;\allowbreak\ 0{,}9647;\allowbreak\ 0{,}01767)$. Como $e^{0}=1$, el futuro conserva peso.
\end{enumerate}
\resultado \textbf{(a) la fila suma $0{,}883$ y los pesos del pasado son menores de lo debido} (habría que renormalizar); \textbf{(b) el token futuro recibe un $1{,}8\,\%$ del peso}: el modelo «ve» el futuro.
\comprobacion la versión (a) renormalizada coincide con la correcta, como predice el \cref{eje:3-renorm}.
\end{solucion}

\begin{solucion}{eje:3-jm}
\datos fila $(0,4,2)$, tercer token prohibido.
\planteamiento calcular $\softmax(0,4,-\infty)$ y $\softmax(0,4,0)$.
\pasos
\begin{enumerate}[label=\arabic*.]
  \item Con $-\infty$: exponenciales $1$; $54{,}60$; $0$. Suma $55{,}60$. Pesos $0{,}01799$; $0{,}9820$; $0$.
  \item Con~0: exponenciales $1$; $54{,}60$; $1$. Suma $56{,}60$. Pesos $0{,}01767$; $0{,}9647$; $0{,}01767$.
  \item Con~0, la posición prohibida recibe peso, y además el resultado depende del valor que tuviera la puntuación original sustituida, lo que es arbitrario.
\end{enumerate}
\resultado \textbf{con $-\infty$: $(0{,}018;\allowbreak\ 0{,}982;\allowbreak\ 0)$; con~0: $(0{,}0177;\allowbreak\ 0{,}9647;\allowbreak\ 0{,}0177)$.} Solo $-\infty$ anula exactamente el peso, porque $e^{-\infty}=0$ mientras que $e^{0}=1$.
\comprobacion la fila con $-\infty$ coincide con la del \cref{ej:mascara}; ambas filas suman~1.
\end{solucion}

\begin{solucion}{eje:3-calculo}
\datos $S$ escalada como en el enunciado; valores $4$, $8$, $12$.
\planteamiento sumar la máscara causal y aplicar la \cref{prop:restringida} fila a fila.
\pasos
\begin{enumerate}[label=\arabic*.]
  \item Con máscara, las entradas por encima de la diagonal pasan a $-\infty$.
  \item Fila~1: solo queda $e^{1}$; pesos $(1,0,0)$.
  \item Fila~2: $e^{0}=1$, $e^{\ln3}=3$; suma~4; pesos $(1/4,3/4,0)$.
  \item Fila~3: $e^{0}=1$, $e^{0}=1$, $e^{\ln2}=2$; suma~4; pesos $(1/4,1/4,1/2)$.
  \item Salidas: $o_1=4$; $o_2=\tfrac14\cdot4+\tfrac34\cdot8=1+6=7$; $o_3=\tfrac14\cdot4+\tfrac14\cdot8+\tfrac12\cdot12=1+2+6=9$.
\end{enumerate}
\resultado
\[
A=\begin{pmatrix}1&0&0\\1/4&3/4&0\\1/4&1/4&1/2\end{pmatrix},\qquad (o_1,o_2,o_3)=(4,7,9).
\]
Sustituir los cincos por cienes \textbf{no cambia nada}: son puntuaciones prohibidas y, por la \cref{prop:restringida}, los pesos no dependen de ellas.
\comprobacion cada fila suma~1 y cada salida está entre el mínimo y el máximo de los valores permitidos: $o_2\in[4,8]$, $o_3\in[4,12]$.
\end{solucion}

\begin{solucion}{eje:3-finito}
\datos fila $(0,4,2)$; se suma $-c$ al tercer elemento; umbral $10^{-6}$.
\planteamiento el peso filtrado es $w(c)=e^{2-c}/(1+e^{4}+e^{2-c})$, decreciente en $c$; se despeja $c$ de $w(c)<10^{-6}$.
\pasos
\begin{enumerate}[label=\arabic*.]
  \item $1+e^{4}=55{,}60$. La condición $w<10^{-6}$ equivale a $e^{2-c}(1-10^{-6})<10^{-6}\cdot55{,}60$, es decir, $e^{2-c}<5{,}560\cdot10^{-5}$ (el factor $1-10^{-6}$ es despreciable).
  \item Tomando logaritmos: $2-c<\ln(5{,}560\cdot10^{-5})=-9{,}797$, luego $c>11{,}80$.
  \item El menor entero es $c=12$. Verificación: $w(12)=e^{-10}/(55{,}60+e^{-10})=4{,}540\cdot10^{-5}/55{,}60=8{,}17\cdot10^{-7}$; $w(11)=e^{-9}/55{,}60=2{,}22\cdot10^{-6}$.
\end{enumerate}
\resultado \textbf{$c=12$}, con un peso filtrado de $8{,}17\cdot10^{-7}$.
\comprobacion el umbral depende de la diferencia entre la puntuación prohibida y las permitidas, que no está acotada a priori; por eso se usa $-\infty$ o el menor número representable, y no un valor fijo pequeño.
\end{solucion}

\begin{solucion}{eje:3-renorm}
\datos $p_j=e^{z_j}/Z$ con $Z=\sum_le^{z_l}$; $\tilde p_j=p_j\mathbb{1}[j\in\mathcal{P}_i]/\sum_{l\in\mathcal{P}_i}p_l$.
\planteamiento sustituir $p_j$ y simplificar $Z$; comparar con la \cref{prop:restringida}.
\pasos
\begin{enumerate}[label=\arabic*.]
  \item Para $j\in\mathcal{P}_i$: $\tilde p_j=\dfrac{e^{z_j}/Z}{\sum_{l\in\mathcal{P}_i}e^{z_l}/Z}=\dfrac{e^{z_j}}{\sum_{l\in\mathcal{P}_i}e^{z_l}}$.
  \item Para $j\notin\mathcal{P}_i$: $\tilde p_j=0$.
  \item Ambas expresiones coinciden con las de la \cref{prop:restringida}; luego $\tilde p=\softmax(z+\mu)$.
  \item Preferencia por la máscara previa: evita una segunda pasada de normalización y, sobre todo, es numéricamente robusta. Si una puntuación prohibida es mucho mayor que las permitidas, los $p_j$ permitidos pueden redondearse a~0 en coma flotante y la renormalización produce $0/0$; restando el máximo \emph{de las permitidas}, la máscara previa nunca lo hace.
\end{enumerate}
\resultado \textbf{$\tilde p=\softmax(z+\mu)$}: anular y renormalizar es matemáticamente equivalente, pero más caro y menos estable.
\comprobacion con $z=(0,4,2)$ y $\mathcal{P}_i=\{1,2\}$ se obtiene $(0{,}01799;\allowbreak\ 0{,}9820;\allowbreak\ 0)$ por ambos caminos (\cref{eje:3-despues}).
\end{solucion}

% ---------------------------------------------------------------------
\section{Tema 4}

\begin{solucion}{eje:4-vf}
\planteamiento contrastar cada afirmación con la \cref{prop:param}, la \cref{prop:equivariancia} y el \cref{cor:paralelo}.
\pasos
\begin{enumerate}[label=(\alph*)]
  \item \textbf{Falsa.} Los parámetros son $4d^{2}$ para todo $h$: cambiar $h$ solo cambia cómo se cortan $\WQ$, $\WK$ y $\WV$ en bloques de $d\times d/h$.
  \item \textbf{Verdadera.} Por la \cref{prop:equivariancia}, $f(PX)=Pf(X)$: cada token recibe la misma salida, en otra fila. Por eso se añade información de posición (módulo~14).
  \item \textbf{Falsa.} Por el \cref{cor:paralelo}, una sola pasada con máscara produce las $n$ predicciones; la máscara es precisamente lo que permite entrenar en paralelo.
\end{enumerate}
\resultado \textbf{(a) falsa; (b) verdadera; (c) falsa.}
\comprobacion para (a), con $d=512$: $4\cdot512^{2}=1\,048\,576$ tanto con $h=8$ (bloques de $64$ columnas) como con $h=16$ (bloques de $32$).
\end{solucion}

\begin{solucion}{eje:4-param}
\datos $d=512$, $h=8$, sin sesgos.
\planteamiento \cref{prop:param}; cada $\WQ^{(r)}$ es $d\times d/h$.
\pasos
\begin{enumerate}[label=\arabic*.]
  \item $d^{2}=262\,144$; $4d^{2}=1\,048\,576$.
  \item $d/h=64$; $\WQ^{(r)}\in\R^{512\times64}$ tiene $32\,768$ parámetros.
  \item $8\cdot32\,768=262\,144=512^{2}$: las ocho juntas, yuxtapuestas por columnas, forman $\WQ\in\R^{512\times512}$.
\end{enumerate}
\resultado \textbf{$1\,048\,576$ parámetros; $32\,768$ por cada $\WQ^{(r)}$.}
\comprobacion coincide con la cifra dada tras la \cref{prop:param} para $d=512$.
\end{solucion}

\begin{solucion}{eje:4-flops}
\datos (a) $n=2048$, $d_k=128$; (b) $d=512$, $n=2048$.
\planteamiento \cref{lem:flops} y \cref{prop:coste}.
\pasos
\begin{enumerate}[label=\arabic*.]
  \item (a) $QK\T$ es $(2048\times128)(128\times2048)$: $2\cdot2048\cdot128\cdot2048=2\cdot4\,194\,304\cdot128=1\,073\,741\,824\approx1{,}074\cdot10^{9}$ FLOP, exactamente $2^{30}$.
  \item $AV$ es $(2048\times2048)(2048\times128)$ y cuesta lo mismo; en total, $2\,147\,483\,648\approx2{,}147\cdot10^{9}$ FLOP.
  \item (b) Proyecciones: $8nd^{2}=8\cdot2048\cdot262\,144=4\,294\,967\,296\approx4{,}295\cdot10^{9}$.
  \item Atención: $4n^{2}d=4\cdot4\,194\,304\cdot512=8\,589\,934\,592\approx8{,}590\cdot10^{9}$.
  \item Cociente: $4n^{2}d/(8nd^{2})=n/(2d)=2048/1024=2$; como $n=2048>2d=1024$, domina la parte cuadrática.
\end{enumerate}
\resultado \textbf{(a) $1{,}074\cdot10^{9}$ FLOP para $QK\T$ y $2{,}147\cdot10^{9}$ con $AV$; (b) $4{,}295\cdot10^{9}$ frente a $8{,}590\cdot10^{9}$: la atención cuesta el doble}, de acuerdo con el criterio $n>2d$.
\comprobacion con potencias de~2: $2\cdot2^{11}\cdot2^{7}\cdot2^{11}=2^{30}$ y $8\cdot2^{11}\cdot2^{18}=2^{32}=4\,294\,967\,296$.
\end{solucion}

\begin{solucion}{eje:4-mem}
\datos fp16, $2$ bytes por número; $n^{2}$ números por cabeza y capa; $1\text{ GB}=10^{9}$ bytes.
\planteamiento memoria $=2n^{2}$ bytes; el coste de la atención es proporcional a $n^{2}$.
\pasos
\begin{enumerate}[label=\arabic*.]
  \item $n=10^{5}$: $2\cdot10^{10}$ bytes $=20$ GB.
  \item $n=10^{6}$: $2\cdot10^{12}$ bytes $=2000$ GB.
  \item $8192/1024=8$; el coste se multiplica por $8^{2}=64$.
\end{enumerate}
\resultado \textbf{$20$ GB; $2000$ GB; factor $64$.}
\comprobacion de $10^{5}$ a $10^{6}$, $n$ se multiplica por~10 y la memoria por $100$: $20\cdot100=2000$, como en la \cref{tab:memoria}.
\end{solucion}

\begin{solucion}{eje:4-jm}
\datos $L=6$, $d=512$, $A=8$, $d_k=d_v=64$, $d_{ff}=2048$.
\planteamiento contar las entradas de cada matriz y sumar.
\pasos
\begin{enumerate}[label=\arabic*.]
  \item (a) $\WQ,\WK,\WV$: $A$ cabezas de $d\times d/A$, es decir, $d^{2}$ cada una; con $\WO$, $4d^{2}=4\cdot262\,144=1\,048\,576$. No depende de $A$, porque $A\cdot d/A=d$.
  \item (b) Red de alimentación: $d\cdot d_{ff}+d_{ff}\cdot d=2\cdot4d^{2}=8d^{2}=2\,097\,152$.
  \item (c) Por capa, $4d^{2}+8d^{2}=12d^{2}=3\,145\,728$; la pila, $12Ld^{2}=6\cdot3\,145\,728=18\,874\,368$.
\end{enumerate}
\resultado \textbf{(a) $1\,048\,576$, independiente de $A$; (b) $2\,097\,152$; (c) $12Ld^{2}=18\,874\,368\approx1{,}89\cdot10^{7}$.}
\comprobacion la atención es $4/12=1/3$ de cada bloque. Los embeddings y la capa de salida, de $\abs{\V}\,d$ parámetros, no están incluidos en $12Ld^{2}$.
\end{solucion}

\begin{solucion}{eje:4-cs336}
\datos $L=12$, $d=768$, $\abs{\V}=50\,257$; $n=1024$ y $n=16\,384$.
\planteamiento sustituir en cada fórmula, multiplicar por $L$ las que son por capa y sumar. Se trabaja en unidades de $10^{9}$ FLOP.
\pasos
\begin{enumerate}[label=\arabic*.]
  \item Auxiliares con $n=1024$: $d^{2}=589\,824$; $n^{2}=1\,048\,576$; $nd^{2}=603\,979\,776$.
  \item Proyecciones ($8nd^{2}$ por capa): $12\cdot8\cdot603\,979\,776=57{,}98\cdot10^{9}$.
  \item Atención ($4n^{2}d$ por capa): $12\cdot4\cdot1\,048\,576\cdot768=38{,}65\cdot10^{9}$.
  \item Red de alimentación ($16nd^{2}$ por capa): $12\cdot16\cdot603\,979\,776=116{,}0\cdot10^{9}$.
  \item Capa de salida: $2\cdot1024\cdot768\cdot50\,257=79{,}05\cdot10^{9}$.
  \item Total: $291{,}6\cdot10^{9}$ (exactamente $291\,648\,307\,200$). Porcentajes: proyecciones $19{,}9\,\%$, atención $13{,}3\,\%$, red de alimentación $39{,}8\,\%$, salida $27{,}1\,\%$.
  \item (b) Con $n=16\,384=16\cdot1024$, los términos lineales en $n$ se multiplican por~16 y la atención por $16^{2}=256$: proyecciones $927{,}7$, atención $9896$, red de alimentación $1855$, salida $1265$ (en $10^{9}$ FLOP).
  \item Total: $13\,944\cdot10^{9}\approx1{,}394\cdot10^{13}$; cociente $13\,944/291{,}6=47{,}8$. Porcentajes: atención $71{,}0\,\%$, red de alimentación $13{,}3\,\%$, salida $9{,}1\,\%$, proyecciones $6{,}7\,\%$.
  \item (c) La atención supera a la red de alimentación cuando $4n^{2}d>16nd^{2}$, es decir, $n>4d=3072$.
\end{enumerate}
\resultado \textbf{con $n=1024$, $2{,}916\cdot10^{11}$ FLOP y domina la red de alimentación ($39{,}8\,\%$); con $n=16\,384$, $1{,}394\cdot10^{13}$ FLOP ($\times47{,}8$) y domina la atención ($71{,}0\,\%$).} La atención supera a la red de alimentación a partir de $n=3072$ y a las proyecciones a partir de $n=2d=1536$.
\comprobacion los porcentajes suman $100{,}1\,\%$ por redondeo en ambos casos. El total crece menos que $256$ veces porque solo una parte del coste es cuadrática.
\end{solucion}

\begin{solucion}{eje:4-rango}
\datos $\WQ,\WK\in\R^{d\times d_k}$ con $d_k<d$; en (c), $d=512$, $h=8$, $d_k=64$.
\planteamiento usar que el rango de un producto no supera el de sus factores y la factorización de rango (módulo~3).
\pasos
\begin{enumerate}[label=\arabic*.]
  \item (a) Para cualesquiera $\WQ,\WK$, $\rg(\WQ\WK\T)\le\rg\WQ\le d_k<d$. La matriz $B=I_d$ tiene rango $d$, luego no se escribe como $\WQ\WK\T$.
  \item (b) Si $\rg B=r\le d_k$, sean $c_1,\dots,c_r$ una base de su espacio de columnas y $C\in\R^{d\times r}$ la matriz que las tiene por columnas. Cada columna de $B$ es combinación de ellas: $B=CR$ con $R\in\R^{r\times d}$. Basta tomar $\WQ=(C\ \ 0)$ y $\WK=(R\T\ \ 0)$, completadas con $d_k-r$ columnas nulas, para que $\WQ\WK\T=CR=B$.
  \item (c) La suma de ocho matrices de rango $\le64$ tiene rango $\le8\cdot64=512=d$, que puede ser completo.
  \item Pero la atención multi-cabeza no aplica una softmax a la suma de las puntuaciones: aplica una softmax distinta a cada cabeza y suma después las salidas (\cref{prop:descomp}). Por la no linealidad de la softmax, $\sum_r\softmax(S_r)\ne\softmax(\sum_rS_r)$ en general.
\end{enumerate}
\resultado \textbf{(a) $I_d$ no es representable; (b) toda forma de rango $\le d_k$ sí; (c) rango $\le512$, pero la suma no describe la multi-cabeza, que combina $h$ distribuciones de atención distintas.}
\comprobacion ejemplo de la no linealidad con dos cabezas de una sola fila y puntuaciones $S_1=(2,0)$, $S_2=(0,0)$: $\softmax(S_1)+\softmax(S_2)=(0{,}8808+0{,}5;\allowbreak\ 0{,}1192+0{,}5)=(1{,}381;\allowbreak\ 0{,}6192)$, mientras que $2\,\softmax(S_1+S_2)=2\,\softmax(2,0)=(1{,}762;\allowbreak\ 0{,}2384)$.
\end{solucion}

% ---------------------------------------------------------------------
\section{Tema 5}

\begin{solucion}{eje:5-capa}
\datos secuencia «el señor Dursley \dots el señor»; la posición final contiene «señor».
\planteamiento aplicar la \cref{prop:una-capa}: los pesos de la posición final solo comparan $x_i$ con cada $x_j$.
\pasos
\begin{enumerate}[label=\arabic*.]
  \item Para predecir «Dursley», la posición final debería atender a la posición de «Dursley», cuyo rasgo distintivo es que \emph{su token anterior} es «señor».
  \item En una capa, la puntuación de esa posición depende de los embeddings de «señor» (consulta) y de «Dursley» y su posición (clave). Una cabeza podría aprender a preferir «Dursley», pero solo como token concreto, no como «el token que siguió a señor».
  \item Si «Dursley» nunca se vio en el entrenamiento, ninguna preferencia aprendida lo favorece; y si el nombre fuera otro, haría falta otra preferencia distinta.
  \item Lo necesario es que una capa anterior haya escrito en la posición de «Dursley» la información «mi token anterior es señor» (cabeza de token anterior).
\end{enumerate}
\resultado \textbf{con una sola capa no puede seleccionarse la posición por su token anterior; hace falta que antes se haya escrito en la posición de «Dursley» la identidad de su predecesor.}
\comprobacion es exactamente la composición K de dos capas descrita tras la \cref{prop:una-capa}.
\end{solucion}

\begin{solucion}{eje:5-previo}
\datos $\alpha_{i,i-1}=e^{\beta^{2}}/(e^{\beta^{2}}+i)$ (\cref{teo:previo}).
\planteamiento sustituir; para (b), la condición es más exigente cuanto mayor es $i$, luego basta imponerla en $i=9$.
\pasos
\begin{enumerate}[label=\arabic*.]
  \item (a) $\beta=2$: $e^{4}=54{,}60$; $\alpha_{5,4}=54{,}60/(54{,}60+5)=54{,}60/59{,}60=0{,}9161$.
  \item (b) $\alpha\ge0{,}99$ equivale a $e^{\beta^{2}}\ge0{,}99\,(e^{\beta^{2}}+i)$, es decir, $0{,}01\,e^{\beta^{2}}\ge0{,}99\,i$, o $e^{\beta^{2}}\ge99\,i$.
  \item Con $i=9$: $e^{\beta^{2}}\ge891$, $\beta^{2}\ge\ln891=6{,}792$, $\beta\ge2{,}606$.
\end{enumerate}
\resultado \textbf{(a) $\alpha_{5,4}=0{,}9161$; (b) $\beta\ge2{,}606$.}
\comprobacion con $\beta^{2}=\ln891$, $\alpha_{9,8}=891/(891+9)=891/900=0{,}99$ exactamente.
\end{solucion}

\begin{solucion}{eje:5-hop}
\datos $\beta=1$; $X=\begin{pmatrix}2&0\\0&2\end{pmatrix}$ (columnas $x_1,x_2$); $\xi^{(0)}=(1;\allowbreak\ 0{,}5)\T$.
\planteamiento $\xi^{(t+1)}=X\softmax(X\T\xi^{(t)})$.
\pasos
\begin{enumerate}[label=\arabic*.]
  \item $X\T\xi^{(0)}=(2,1)$. Softmax: $e^{2}/(e^{2}+e)=1/(1+e^{-1})=0{,}7311$ y $0{,}2689$.
  \item $\xi^{(1)}=0{,}7311\,(2,0)+0{,}2689\,(0,2)=(1{,}462;\allowbreak\ 0{,}5379)$.
  \item $X\T\xi^{(1)}=(2{,}924;\allowbreak\ 1{,}076)$; diferencia $1{,}848$; pesos $1/(1+e^{-1{,}848})=0{,}8640$ y $0{,}1360$.
  \item $\xi^{(2)}=0{,}8640\,(2,0)+0{,}1360\,(0,2)=(1{,}728;\allowbreak\ 0{,}2721)$.
\end{enumerate}
\resultado \textbf{$\xi^{(1)}=(1{,}462;\allowbreak\ 0{,}5379)$, $\xi^{(2)}=(1{,}728;\allowbreak\ 0{,}2721)$}: el estado se acerca al patrón $x_1=(2,0)$, el más parecido al estímulo inicial.
\comprobacion una tercera actualización da pesos $(0{,}9484;\allowbreak\ 0{,}0516)$ y $\xi^{(3)}=(1{,}897;\allowbreak\ 0{,}1032)$, aún más cerca de $x_1$; cada estado está en el segmento que une $x_1$ y $x_2$ (\cref{prop:convexa}).
\end{solucion}

\begin{solucion}{eje:5-induccion}
\datos secuencia «$\langle s\rangle$ el gato come pescado y luego el gato duerme», posiciones $0$ a~$9$; $e_t$ \emph{one-hot} de token, $p_i$ \emph{one-hot} de posición; $\beta=10$; máscara causal; sin escalado adicional.
\planteamiento capa~1: cabeza de token anterior (\cref{teo:previo}); capa~2: consulta con el token actual y claves con la salida de la capa~1 (composición K).
\pasos
\begin{enumerate}[label=\arabic*.]
  \item Capa~1: $q_i=\beta p_i$, $k_j=\beta p_{j+1}$, $v_j=e_{x_j}$. Por el \cref{teo:previo}, $\alpha_{i,i-1}=e^{100}/(e^{100}+i)\approx1$, y la salida es $h_i\approx e_{x_{i-1}}$ para $i\ge1$ (en $i=0$, $h_0=e_{\langle s\rangle}$).
  \item Capa~2: $q_i=\beta e_{x_i}$, $k_j=\beta h_j$, $v_j=e_{x_j}$. La puntuación $q_i\cdot k_j\approx\beta^{2}=100$ si $x_{j-1}=x_i$ y $\approx0$ en otro caso.
  \item Posición~7 («el»): entre $j\le7$, solo $j=2$ cumple $x_{j-1}=x_1=$ «el» (la propia $j=8$ está prohibida por la máscara). Peso $\approx1$ en $j=2$; valor $e_{\text{gato}}$: predice \textbf{«gato»}.
  \item Posición~8 («gato»): solo $j=3$ cumple $x_{2}=$ «gato». Valor $e_{\text{come}}$: predice \textbf{«come»}.
  \item Posición~9 («duerme»): ninguna $j\le9$ tiene «duerme» como token anterior; todas las puntuaciones valen~0 y la atención es uniforme, $1/10$ por posición: no se copia nada útil.
\end{enumerate}
\resultado \textbf{capa~1: $Q=\beta\cdot$posición, $K=\beta\cdot$posición desplazada, $V=$ token; capa~2: $Q=\beta\cdot$token actual, $K=\beta\cdot$salida de la capa~1, $V=$ token.} Predice «gato» tras el segundo «el» y «come» tras el segundo «gato»; en la posición~9 no hay patrón previo.
\comprobacion una simulación en NumPy da pesos $1{,}0000$ en la posición~2 (para $i=7$) y en la~3 (para $i=8$), y peso máximo $0{,}1=1/10$ en $i=9$.
\end{solucion}

\begin{solucion}{eje:5-jac}
\datos $o(\xi)=\sum_lp_lx_l$, $p_l=e^{\beta x_l\T\xi}/Z$, $Z=\sum_me^{\beta x_m\T\xi}$.
\planteamiento derivar $p_l$ respecto de $\xi$ con la regla del cociente (módulo~4) y sustituir.
\pasos
\begin{enumerate}[label=\arabic*.]
  \item $\nabla_\xi e^{\beta x_l\T\xi}=\beta e^{\beta x_l\T\xi}x_l$ y $\nabla_\xi Z=\beta\sum_me^{\beta x_m\T\xi}x_m=\beta Z\,o$.
  \item Regla del cociente: $\nabla_\xi p_l=\bigl(\beta e^{\beta x_l\T\xi}x_lZ-e^{\beta x_l\T\xi}\beta Zo\bigr)/Z^{2}=\beta p_l(x_l-o)$.
  \item $\dfrac{\partial o}{\partial\xi}=\sum_lx_l(\nabla_\xi p_l)\T=\beta\sum_lp_lx_l(x_l-o)\T=\beta\Bigl(\sum_lp_lx_lx_l\T-o\,o\T\Bigr)$, usando $\sum_lp_lx_l=o$.
  \item La matriz entre paréntesis es $\E_p[xx\T]-\E_p[x]\E_p[x]\T=\Cov_p(x)$.
  \item (b) Patrones $0$ y $2$, $\xi=\tfrac12\ln3$: puntuaciones $0$ y $\ln3$; pesos $1/4$ y $3/4$; $o=\tfrac34\cdot2=1{,}5$; $\E_p[x^{2}]=\tfrac34\cdot4=3$; varianza $3-1{,}5^{2}=0{,}75$.
  \item (c) Si $p$ se concentra en un patrón, $\Cov_p(x)\to0$ y la jacobiana se anula.
\end{enumerate}
\resultado \textbf{$\partial o/\partial\xi=\beta\Cov_p(x)$; en el ejemplo, $do/d\xi=0{,}75$; con $p$ concentrada, la jacobiana tiende a~0.}
\comprobacion la diferencia finita $[o(\xi+10^{-6})-o(\xi-10^{-6})]/(2\cdot10^{-6})$ da $0{,}7500$. El apartado (c) es la saturación de la \cref{prop:saturacion} vista desde la salida: un punto fijo bien separado es estable porque pequeñas perturbaciones del estado apenas cambian la salida.
\end{solucion}

\begin{solucion}{eje:5-recup}
\datos $x_1,\dots,x_N$ ortonormales; $\xi=x_k$; $\xi^{\text{nuevo}}=X\softmax(\beta X\T\xi)$.
\planteamiento calcular los pesos y la norma del error con el teorema de Pitágoras.
\pasos
\begin{enumerate}[label=\arabic*.]
  \item Por ortonormalidad, $X\T x_k=e_k$: puntuación~1 en $k$ y~0 en el resto.
  \item Pesos: $p_k=e^{\beta}/(e^{\beta}+N-1)$ y $p_l=w:=1/(e^{\beta}+N-1)$ para $l\ne k$. Además $1-p_k=(N-1)w$.
  \item $\xi^{\text{nuevo}}-x_k=-(1-p_k)x_k+\sum_{l\ne k}w\,x_l$.
  \item Como los $x_l$ son ortonormales: $\norm{\xi^{\text{nuevo}}-x_k}^{2}=(1-p_k)^{2}+(N-1)w^{2}=(N-1)^{2}w^{2}+(N-1)w^{2}=N(N-1)w^{2}$.
  \item Raíz: $\norm{\xi^{\text{nuevo}}-x_k}=\sqrt{N(N-1)}\,w$.
  \item Con $N=4$, $\beta=5$: $e^{5}=148{,}4$; $w=1/151{,}4=0{,}006605$; $\sqrt{12}=3{,}464$; error $3{,}464\cdot0{,}006605=0{,}02288$.
\end{enumerate}
\resultado \textbf{$\norm{\xi^{\text{nuevo}}-x_k}=\sqrt{N(N-1)}/(e^{\beta}+N-1)$; con $N=4$ y $\beta=5$, $0{,}02288$.}
\comprobacion el error decrece como $e^{-\beta}$: es el error de recuperación exponencialmente pequeño de~\cite[§2]{ramsauer} en su caso más sencillo. El cálculo directo con la matriz identidad de $4\times4$ da el mismo valor.
\end{solucion}
